%=======================================================================
%  CHAPTER 2 -- RELATED WORK
%=======================================================================
\chapter{Background and Related Work}
\label{ch:related}

%-----------------------------------------------------------------------
\section{Anomaly Detection in Multivariate Time Series}
\label{sec:rw-ad}

The classical approach to fault detection rests on statistical process control and its multivariate extensions, principally principal component analysis and its dynamic and kernel variants. Chiang et al.~\cite{chiang2001faultdetection} formalise this lineage as the foundation of data-driven fault diagnosis for industrial systems. In practice such methods frequently reduce to evaluating univariate thresholds, because where the relational structure between variables carries little information the multivariate machinery has nothing to exploit that a threshold does not.

That this remains the operational default is documented empirically, not assumed here. Yin et al.~\cite{yin2012comparison} establish classical multivariate statistical methods as the standard data-driven baseline for the Tennessee Eastman Process, the same dataset used in Chapter~\ref{ch:experimental_setup}. A parameter-free statistical baseline is therefore not a straw man constructed for contrast; it is the method a practitioner would already have deployed, and against which any structural or deep model must justify its additional cost.

This thesis codifies that baseline as a parameter-free mean absolute $z$-score over the monitored signal, evaluated under the same threshold rule as every other configuration. Including it is what allows the evaluation to distinguish a pipeline that fails to learn from one that learns structure it did not need. When Chapter~\ref{ch:discussion} reports that this control scores 0.9286 on the Industrial domain against 0.9151 for the best causal configuration, the finding is not that structural machinery is redundant in principle, but that on this domain, under this detector, it bought nothing the classical mechanism did not already provide.

The remainder of this chapter turns to what has been proposed since: reconstruction-based deep models in Section~\ref{sec:rel-neural}, the evaluation practices under which they are compared in Section~\ref{sec:rel-evaluation}, and the causal methods that are the subject of this thesis from Section~\ref{sec:rw-cd} onward.

\section{Deep Learning Baselines in Time-Series Anomaly Detection}
\label{sec:rel-neural}

The taxonomy of anomaly detection was set out by Chandola et al.~\cite{chandola2009survey} and brought up to date for time series by Bl\'azquez-Garc\'ia et al.~\cite{blazquezgarcia2021review}. Contemporary practice is dominated by reconstruction-based deep models, among them the LSTM encoder-decoder of Malhotra et al.~\cite{malhotra2016lstmad}, the stochastic recurrent model of Su et al.~\cite{su2019omnianomaly}, and the adversarially trained USAD of Audibert et al.~\cite{audibert2020usad}.

Those specialised architectures were not selected as the baselines for this study, and the reason is methodological, not dismissive. It follows from a requirement the comparison has to meet, which is worth stating before the choice that meets it.

A baseline that is tuned less thoroughly than the proposed method will lose, and the result will report the depth of the tuning, not the merit of the model. Section~\ref{sec:rel-evaluation} returns to this as a documented failure of the field, and it is a failure this project committed and had to correct. Avoiding it requires giving every model the same tuning effort, which in practice means one hyperparameter grid searched identically for each.

The specialised architectures make that impossible. Their hyperparameters are not merely different values, they are different quantities: USAD is governed by the weighting between its two adversarially trained decoders, OmniAnomaly by the weight on a variational term in its loss, and neither has any counterpart in the other or in a plain autoencoder. There is no grid that can be searched identically over all of them, because there is no shared axis to vary. One could tune each model on its own terms, but then the models would have received different amounts of search and the comparison would be back where it started.

The baselines here are instead an LSTM autoencoder~\cite{hochreiter1997lstm}, a temporal convolutional network~\cite{bai2018tcn} and a Transformer encoder~\cite{vaswani2017attention}. These span the three mainstream sequence-modelling paradigms (recurrence, convolution and attention), so the comparison still covers the space of contemporary architectures. What makes them usable is that all three reconstruct a window under the same objective and train through the same loop, differing only in how the encoder is built. Their hyperparameters are therefore the same three quantities in every case: hidden width, depth and learning rate. One grid over those three is meaningful for all of them, which is what allows the identical twelve-configuration search described in Chapter~\ref{ch:experimental_setup}, and the argument of Section~\ref{sec:rel-evaluation} depends on their having received it.

Thresholding is the second point of departure. Hundman et al.~\cite{hundman2018telemanom} popularised dynamic, nonparametric thresholds that adapt to the local behaviour of the score series. This thesis uses a fixed 95th percentile of the fault-free score distribution instead. Both rules are label-blind and neither is more principled in the abstract, but only the fixed rule holds the false positive budget identical across configurations, and without that the 27 causal configurations and the baselines could not be placed on a common axis.

The most directly relevant work is the Graph Deviation Network of Deng and Hooi~\cite{deng2021gdn}, which learns a graph over sensors and uses it to detect anomalies. GDN does not claim that the learned graph represents physical causation; the structure is a modelling device. That makes it the sharpest available point of comparison, because it isolates the question this thesis asks: what does the causal claim add beyond structure?

Part of the answer is uncomfortable and is developed in Chapter~\ref{ch:discussion}. On the evidence gathered here, detection performance cannot distinguish a causal graph from a merely structural one: a graph carrying no observed-parent edges at all detected competitively, and two graphs built on incompatible premises produced indistinguishable scores. Judged by detection alone, the weaker claim GDN makes is the better-supported one. The causal framing earns its keep on a different axis: it yields a root-cause attribution that can be scored against a documented ground truth, which a structural graph learned for detection does not attempt and detection metrics cannot assess. That distinction, not any difference in AUC-ROC, is what the causal label buys.

\section{Evaluation Practice in Time-Series Anomaly Detection}
\label{sec:rel-evaluation}

Progress in time-series anomaly detection is frequently illusory, driven by evaluation practice rather than by algorithmic improvement~\cite{wu2023flawed}. Large-scale empirical review supports the weaker but still awkward claim that architectural elaboration does not reliably buy accuracy: Schmidl et al.~\cite{schmidl2022anomaly} re-implemented 71 algorithms and ran them over 976 datasets, and report no clear winner between algorithm families, with every family capable of being effective and the best cost-benefit ratio in the study belonging to a wavelet-based method and not to any deep architecture. A field in which elaborate and simple methods cannot be separated on accuracy is one where the comparison protocol, not the model, decides what looks best. The most consequential of those protocol choices is point adjustment, under which an entire anomalous segment is credited as detected whenever any single point inside it is flagged. Kim et al.~\cite{kim2022rigorous} show that this inflates reported scores so severely that even a random anomaly score can be made to look state of the art. It is not a failure this thesis has to correct for, because no metric here is point-adjusted: every figure reported is computed point-wise against the per-sample labels, so a detector is credited only for the samples it actually flags. Four further failures are not disposed of so easily, and the methodology is structured to measure them instead of merely avoiding them.

The first is fitting the decision threshold to the evaluation labels, which selects the operating point with knowledge of the answers. The pipeline here uses a label-blind 95th-percentile rule instead, and the rule is enforced automatically instead of by convention: every configuration is audited for any step that would choose the operating point with knowledge of the evaluation labels, and the build fails if one is found.

The second is evaluating an optimised proposal against an untuned baseline. This error was committed during the early stages of this project, which evaluated the causal configurations against neural architectures left at default settings. Correcting it, by applying the identical twelve-configuration search to every architecture on every domain, moved seven of the nine architecture-domain cells off their default configuration. The correction inverted no domain winner, and it reduced the causal margin on Healthcare from 0.1299 against the untuned baseline to 0.1037 against the tuned one. A comparison that omits baseline tuning reports the depth of its own hyperparameter search and not the merit of its algorithm.

The third is the omission of a parameter-free control. Introducing a trivial mean absolute $z$-score baseline changed the conclusion available on the Industrial domain instead of confirming it. The control scores 0.9286 against 0.9151 for the best causal configuration, so the finding is not that causal methods underperformed deep learning there, but that neither method's structure performed work a threshold could not. A comparison omitting the control would have reported the first conclusion, which is materially different and, on this evidence, wrong.

The fourth is treating a numerical difference as a demonstrated one. As Dem\v{s}ar~\cite{demsar2006statistical} sets out, a reported performance gap requires a formal test. Applying DeLong's test across the evaluation grid revealed a structural obstacle, not a result: 54 of the 108 within-domain comparisons score different observations and admit no paired test at all. Every one of those 54 involves PCMCI against one of the other two algorithms, which locates the cost precisely. Treating the subsampling rate as a property of each algorithm's workflow instead of as a pipeline constant is defensible, and it forecloses cross-algorithm significance testing entirely, a consequence worth stating in the same breath as the choice.

\section{Causal Discovery from Observational Data}
\label{sec:rw-cd}

Observational causal discovery is surveyed by Glymour et al.~\cite{glymour2019review} and, more critically, by Vowels et al.~\cite{vowels2022dags}. This thesis structures its evaluation around the three principal families that taxonomy identifies, taking one representative from each, so that a result holding across all three is a property of causal discovery and not of one estimator.

The first family is constraint-based discovery, originating with the PC and FCI algorithms~\cite{spirtes2000causation}, which recover structure through conditional independence testing. Runge et al.~\cite{runge2019pcmci,runge2020pcmciplus} extended the paradigm to time series, where autocorrelation inflates apparent dependence, and the approach is implemented in the Tigramite framework~\cite{runge2023tigramite}. PCMCI is taken as the representative because the standard method degrades sharply under exactly the conditions these domains present. On synthetic data constructed to contain no cross-variable edges at all, so that every reported edge is a false positive by construction, PC links 0.4982 of all variable pairs at an autocorrelation coefficient of 0.95, against 0.0429 at zero, while PCMCI's rate is flat across the whole sweep. That flatness, not any particular level, is what the selection rests on: the level itself depends on the false-discovery correction, which this pipeline does not apply. The measurement, and that distinction, are reported in Section~\ref{sec:pcpcmci}.

The second family is score-based discovery, which searches the space of Markov equivalence classes to optimise a penalised likelihood. Chickering~\cite{chickering2002ges} formalised Greedy Equivalence Search, later scaled to high dimensions by Ramsey et al.~\cite{ramsey2017fges}. Because it scores equivalence classes instead of individual directed graphs, GES returns a completed partially directed acyclic graph, in which some edges remain unoriented. As set out in Chapter~\ref{ch:background}, that property obliges the pipeline to apply a strict extraction rule and retain only unambiguously directed edges, since an unoriented edge does not license a parent-child regression.

The third family comprises functional causal models. The lineage begins with LiNGAM~\cite{shimizu2006lingam}, which obtains identifiability from non-Gaussian noise under linear mechanisms, and continues through nonlinear additive noise models~\cite{hoyer2008anm,buhlmann2014cam}. Maeda and Shimizu~\cite{maeda2021camuv} developed CAM-UV, which also flags pairs of variables it attributes to an unobserved common cause, and which this thesis runs through the lingam package~\cite{ikeuchi2023lingam}. Its identifiability derives from the nonlinearity of the causal mechanisms under additive noise and not from the non-Gaussianity assumption central to the earlier linear variants, a distinction that matters because the two are easily conflated within the same lineage.

Continuous-optimisation approaches targeting independent observations, such as NOTEARS~\cite{zheng2018notears}, fall outside this three-family taxonomy and are excluded on that basis and not on any claim about their capability. The selection principle was one representative per family, and NOTEARS belongs to none of the three.

The software ecosystem is fragmented in a way that shaped the implementation. The causal-learn library~\cite{zheng2024causallearn} supplies GES, and its Table 1 organises the field into the families used above, but it contains no time-series constraint-based method: the constraint-based algorithms it ships assume independent observations. CAM-UV is run from the lingam package and PCMCI from Tigramite, so no single library could supply all three and a unified wrapper across three packages was required, with the verification that entails.

These three families make incompatible assumptions about what makes structure recoverable: conditional independence, penalised likelihood over equivalence classes, and functional form. Evaluating them under one fixed downstream detector is what turns that incompatibility from a theoretical observation into a measurable difference.

\section{Causality for Anomaly Detection and Root-Cause Analysis}
\label{sec:rw-causal-ad}

The nearest neighbours to this work sit in three literatures. Root-cause analysis for microservice systems has produced a sequence of causal-discovery-based methods, among them RCD~\cite{ikram2022rcd}, CIRCA~\cite{li2022circa} and BARO~\cite{pham2024baro}, together with a critical assessment of how far the field has actually progressed~\cite{pham2024howfar}. Fault detection and diagnosis in the process industries has a longer history, surveyed by Chiang et al.~\cite{chiang2001faultdetection} and compared empirically by Yin et al.~\cite{yin2012comparison} on the Tennessee Eastman benchmark used here. And a separate line derives detectors from Granger causality~\cite{granger1969causality}, which tests predictive rather than structural precedence.

Across these literatures a recurring methodological pattern makes reported gains difficult to interpret. A paper typically proposes a causal discovery method, pairs it with a downstream detector designed alongside it, evaluates on a dataset chosen by the authors, and scores the result under a thresholding rule specific to the paper. When the resulting pipeline outperforms an off-the-shelf recurrent baseline, the improvement cannot be attributed to the causal structure, because the causal structure is not the only thing that changed. The detector, the threshold and the evaluation protocol changed with it, and any of them could carry the difference.

This thesis addresses that by fixing the downstream machinery instead of the inputs. No claim of complete isolation is made, and the limits are stated where they arise: as set out in Chapter~\ref{ch:methodology}, the subsampling rate and the variable inclusion cap are properties of each algorithm's own workflow and are not equalised, because equalising them is infeasible in one direction and changes what each algorithm is in the other. What is held constant is everything downstream of the discovered graph (the recursive least squares filter, the coefficient drift score, the fixed 95th-percentile threshold and the attribution depth) identically across all 27 configurations, and verified programmatically on every build rather than asserted.

The contribution is therefore narrower than a claim of priority and more useful than one. Within this evaluation a performance difference between two configurations can be attributed to the causal discovery algorithm and its interaction with the regression family, because nothing else downstream was permitted to move. Where that attribution is not available, across algorithms that score different observations, the thesis reports the fact, not the difference.

\section{Explainability for Anomaly Detection}
\label{sec:rw-xai}

Explainability for time-series models is comparatively under-served relative to vision and language, as the review of Theissler et al.~\cite{theissler2022xaitimeseries} sets out. The prevailing approach is post-hoc attribution, applying a method such as SHAP~\cite{lundberg2017shap}, LIME~\cite{ribeiro2016lime} or integrated gradients~\cite{sundararajan2017integrated} to a model that has already been trained and is otherwise opaque. Some of this
work is specific to sequences: TimeSHAP~\cite{bento2021timeshap} adapts Shapley
values to recurrent models by perturbing whole subsequences, and reports which
events and which time steps carry the attribution as well as which features.
That is the closest post-hoc analogue to what is attempted here, and the
difference is where the explanation comes from. TimeSHAP explains a trained
model after the fact. The attribution in this work is read off the structure the
detector already uses to score, so no second model is fitted to produce it.

The causal attribution used here is of a different kind, and the distinction is the one Rudin~\cite{rudin2019stop} draws in arguing that high-stakes settings call for models that are interpretable by construction rather than explained after the fact. The attribution in this pipeline is a coefficient vector indexed by the causal parents of each target, so the explanation is the model itself and not an approximation of it. Applying SHAP to an autoencoder's reconstruction error is, by contrast, an approximation of a model that has no representation of one variable causing another. SHAP is run over the neural baselines in this thesis precisely so that this remains a comparison instead of an assertion.

The more consequential problem in this literature is evaluation. As Doshi-Velez and Kim~\cite{doshivelez2017rigorous} argue, explanation quality is usually assessed by asking whether an explanation looks reasonable to a reader, for want of anything better. Adebayo et al.~\cite{adebayo2018sanity} show how weak that standard is, demonstrating that widely used saliency methods can fail elementary sanity checks while producing visually convincing output. The evaluation in this thesis avoids the difficulty instead of solving it: the injected attacks in the robotics dataset specify which physical subsystem was perturbed, so an attribution can be scored against a documented answer rather than a judgement.

That advantage produces a result which complicates Rudin's position instead of confirming it, and the complication is reported in Chapter~\ref{ch:discussion} and not smoothed over here. The intrinsic attribution is correct in 9 of 9 configurations for the LED attack and 8 of 9 for the joint attack. It is also correct in six configurations whose graphs contain no LED channel at all, and therefore cannot represent the relation that would explain the answer. An interpretable-by-construction model arrived at the right explanation from a structure that does not contain it, which is not what interpretability by construction is usually taken to promise.

\section{Summary and Positioning}
\label{sec:rw-gap}

The literature reviewed in this chapter leaves three gaps in how causal anomaly detection is evaluated. This thesis is positioned against those three, and the correspondence to its contributions is direct.

First, published comparisons routinely vary the causal discovery method, the downstream detector, the dataset and the thresholding rule at the same time. A reported gain therefore cannot be attributed to the causal structure, because the structure is not the only thing that changed. Contribution C1 addresses this by fixing everything downstream of the discovered graph (the filter, the score, the threshold rule and the attribution depth) and by enforcing that uniformity through automated checks on every build instead of by assertion. What remains variable, the subsampling rate and the variable inclusion cap, is declared as a property of each algorithm's workflow and its cost is quantified in Chapter~\ref{ch:discussion}, so the claim made here is that the downstream machinery is controlled, and not that the comparison is fully isolated.

Second, the absence of a parameter-free reference point leaves claims of algorithmic superiority unanchored. Where a proposed causal method outperforms a neural baseline, the literature rarely establishes whether a trivial statistical threshold would have outperformed both. Contributions C2 and C3 follow from supplying that reference point. C2 uses it as an instrument, not just as a control, since the score of the trivial baseline indexes how much of a domain's fault signature is pure amplitude and the causal advantage falls monotonically as it rises; C3 places the tuned neural architectures on the same axis under the same label-blind protocol, so that all three families are read against one another instead of each against its own convention. The accompanying negative result is reported where it arises but is not claimed as a contribution: detection performance cannot establish that a discovered graph is physically meaningful, because a graph carrying no observed-parent edges detected competitively and graphs built on incompatible premises produced indistinguishable scores.

Third, explanation quality is predominantly assessed by whether an explanation appears plausible, for want of a documented answer to compare against. The explainability half of Contribution C2 exploits a dataset in which the injected faults name the perturbed subsystem, so attribution can be scored rather than left to judgement, and C3 sets that intrinsic attribution against post-hoc SHAP over the neural baselines on the same faults. The result complicates the interpretability literature, and does not confirm it: a model interpretable by construction produced the correct explanation in configurations whose graphs do not contain the relation that would account for it.

Two further findings do not answer a gap in this literature and are not presented as contributions. The first bounds the behaviour of CAM-UV empirically on real data, replacing an assumption about its latent-variable handling with measured saturation and starvation regimes. The second identifies a consequence of this study's own design: that treating the subsampling rate as an algorithm property forecloses cross-algorithm significance testing. Both are findings about the methods and about this evaluation and not about what the field has failed to do, and both are reported where they arise.
